The pushforward turns an explanation over the bytes into one over the features: it explains a model in a vocabulary
chosen by the analyst, here the EMBER features, even when the model never computes them. We ask it the three questions of
Section~\ref{sec:XML:Pullback}, in the same order. Is the transferred explanation faithful
(Section~\ref{sec:XML:PushforwardFaithfulness})? Does it recover a known answer, the planted trigger
(Section~\ref{sec:XML:PushforwardTrigger})? And what does it let a defender do that was not possible before? For the
pushforward, the answer to the last question is an audit of a raw-byte detector across a corpus
(Section~\ref{sec:XML:Audit}). The data, models, and deletion operator are those of Section~\ref{sec:XML:Pullback}.

\subsection{Faithfulness}
\label{sec:XML:PushforwardFaithfulness}

A pushed-forward explanation is useful only if the features it ranks highest are the features the decision rests on.
We test this in the feature space, by deleting the top-ranked features, and compare the pushforward with the explainers
that work there directly. In short, explanations pushed forward from the bytes are far more faithful than random and
than a naive translation, but reach only about a third of TreeSHAP's faithfulness. The loss is not in the pushforward
but in the resolution of the explanations we push forward: the finer the source, the more of TreeSHAP's faithfulness
the pushforward attains (Table~\ref{tab:XML:Resolution}).

We push forward the direct problem-space explanations of Section~\ref{sec:XML:PullbackFaithfulness}, occlusion at 64 to
4{,}096 windows and KernelSHAP over 256 windows, each through $P$ coarsened to its windows. The resulting feature
rankings are scored by a feature-space deletion curve: we delete the top-ranked features, re-score the model, and
compare the area under the curve with those of TreeSHAP, KernelSHAP, random orders, and a naive translation that
credits each feature with the plain sum of the windows it reads (no $D^{\dagger}$). A deleted feature takes its value
in the fully ablated file, the reference that matches the problem-space explanations (deleting every feature reproduces
the end point of the curves of Section~\ref{sec:XML:PullbackFaithfulness} exactly), or, as a robustness check, the
interventional expectation over the background, the reference that matches TreeSHAP. We summarize a method by the ratio
$R$ of its area to TreeSHAP's, and read $R \geq 0.8$ as faithful, 0.5--0.8 as partial, and below 0.5 as weakened.

Every pushed-forward explanation is far above random (97--100 of 100 files, $p \approx 10^{-18}$) and above the naive
translation (57--70 of 100 files, $p = 2.6 \cdot 10^{-6}$ to $0.023$), so the normalization by $D^\dagger$ matters. Yet
they reach only $R = 0.31$--$0.40$ of TreeSHAP's area (0.25--0.29 under the background reference), which is weak. The
pushforward itself is not the cause: TreeSHAP's own explanation, pulled back to the bytes and pushed forward again,
keeps $R = 0.80$. To locate the loss, we express TreeSHAP's information at $K$ windows (pull it back, sum it over the
windows, and push it forward through $P$ coarsened to the same windows) and measure what survives.

\begin{table*}[t]
  \centering
  \small
  \begin{tabular}{lcccccccc}
    \toprule
    Resolution (windows) & 64 & 256 & 1{,}024 & 4{,}096 & 16{,}384 & 65{,}536 & 262{,}144 & bytes \\
    \midrule
    $R$ attainable (TreeSHAP's information) & 0.34 & 0.37 & 0.43 & 0.50 & 0.61 & 0.69 & 0.75 & 0.80 \\
    $R$ of occlusion & 0.31 & 0.32 & 0.34 & 0.37 & -- & -- & -- & -- \\
    \bottomrule
  \end{tabular}
  \caption{What survives the pushforward at each resolution (100 files, path TreeSHAP, the matched reference). $R$ is a
  method's deletion area relative to TreeSHAP's. The attainable $R$ expresses TreeSHAP's own explanation at that
  resolution and pushes it forward; occlusion cannot be computed beyond 4{,}096 windows at an affordable cost.
  KernelSHAP over 256 windows reaches 0.40. Faithfulness rises with the resolution of the source, and occlusion stays
  close to what its resolution allows.}
  \label{tab:XML:Resolution}
\end{table*}

Table~\ref{tab:XML:Resolution} shows that what a $K$-window explanation can attain rises at every step, for every
explainer and both references, from $R = 0.34$ at 64 windows to $0.80$ at the bytes. The coarser the windows, the more
features share each one, and the less of a window's attribution the pushforward can return to the feature it came from
(Section~\ref{sec:XML:Constructing}). Because this attainable $R$ uses the feature-space explainer's own information,
its byte-level end point is partly circular; the evidence is the trend across resolutions. At 256 windows even
TreeSHAP's information reaches only $R = 0.37$, and occlusion (0.32) and KernelSHAP over windows (0.40) are near that
bound; occlusion's share of the bound falls from 0.91 at 64 windows to 0.74 at 4{,}096, where it saturates, as in
Section~\ref{sec:XML:PullbackFaithfulness}. The pushforward is thus faithful up to the resolution of its source: the
direct problem-space explainers are limited by their resolution, not by the transfer.

The test above pushes forward explanations of the GBDT, which lets us compare with TreeSHAP; the audit of
Section~\ref{sec:XML:Audit} instead pushes MalConv's Integrated Gradients forward, so we also test the pushforward on
MalConv, where it is used. MalConv never computes the features, so its features cannot be deleted; instead we pull
each feature explanation back through $P$, delete its top bytes with the operator of
Section~\ref{sec:XML:PullbackFaithfulness}, and let MalConv score the result. The realistic alternative is a surrogate:
a GBDT trained on the EMBER features to predict MalConv's log-odds (on test files outside those explained;
$R^2 = 0.95$ on them), explained by TreeSHAP. On the 100 files, the pushed-forward explanation removes about twice as
much of MalConv's evidence as the surrogate's (median area 6.25 against 3.28 log-odds, larger on 99 of 100 files,
$p \approx 10^{-18}$), and 2.7 times as much as random feature explanations (2.33); with random instead of zero fill, it removes 2.0 times
as much as the surrogate (99 of 99 files). Random feature explanations are
the fair floor here: pulled back, any feature explanation lands on bytes the extractor reads, which MalConv also relies
on, so they already remove far more than random windows (0.61). The surrogate imitates MalConv's outputs closely, but
its explanation describes the surrogate; the pushforward describes MalConv, in the surrogate's vocabulary. The
pushforward also beats Integrated Gradients summed over 256 windows (4.29, larger on 83 of 100 files), which names no
features. Integrated Gradients itself, per byte, removes far more (33.4), but not because it is a better summary:
deleting its top 1\% of bytes already removes 19 log-odds, three times as much as deleting the whole file (6.0), with
either fill, the fragility of MalConv noted in Section~\ref{sec:XML:Comparison}. Any summary into features gives up this byte-level
precision, and among summaries into features, the pushforward is the faithful one.

\subsection{A Planted Trigger}
\label{sec:XML:PushforwardTrigger}

The planted trigger of Section~\ref{sec:XML:PullbackTrigger} gives the pushforward a known answer as well, and in
doing so gives feature-space explainers a ground truth they otherwise lack. On a triggered file, the pushforward of the
trigger's bytes weights each feature by how much of it the trigger computes, a property of $P$ alone. Its support, the
features the trigger touches, does not depend on the model, and a faithful explanation of the trigger's effect must
concentrate there. We score TreeSHAP and KernelSHAP of the poisoned GBDT by the share of their attribution change,
$a(x_T) - a(x)$, that falls on this support, and by whether their top-ranked features lie in it. The weights themselves
are not a target: they cannot say which of the touched features the model uses, and even an ideal explainer of a tree
model moves some credit onto untouched features through interactions.

Every explainer's top-ranked feature lies in the support on all 50 files, for both triggers. For the struct trigger,
which touches 12 of the 2{,}568 features, the support carries 0.70 of the attribution change for both TreeSHAP variants
and 0.50 for KernelSHAP, against 0.005 for a random attribution. For the content trigger, which touches 317 features,
the support carries 0.67 and 0.51 against 0.12, and half of each explainer's ten top-ranked features lie in it.
KernelSHAP spreads more of its change outside the support, as it spreads more of its pullback outside the trigger
(Section~\ref{sec:XML:PullbackTrigger}), and it needs its full budget: with 512 coalitions instead of 32{,}768 it puts
0.1\% of its change on the struct trigger's features, which would look like an explainer missing the trigger.

\subsection{Application: Auditing a Raw-Byte Detector}
\label{sec:XML:Audit}

A detector that reads the raw bytes directly, such as MalConv or the deep models that deployed stacks run over full
file contents~\cite{microsoft2019insideout}, can be explained file by file with Integrated Gradients, but it has no
vocabulary in which its reliance recurs across files: a byte offset means something different in each file. Auditing it
for reliance on content that an attacker could change cheaply, without affecting what the program does, is therefore a
corpus-level question~\cite{arp2022and}, and per-file byte explanations cannot be aggregated to answer it. The
pushforward supplies the shared vocabulary. It sends each file's byte explanation into the named features of a feature
extractor, which are the same for every file, splitting the credit of shared bytes among the features that read them.
This is the global counterpart of the pullback's local comparison (Section~\ref{sec:XML:Comparison}), on the same two
detectors.

\paragraph{Threat model.} An attacker wants a malicious file that MalConv detects to be classified benign. The attacker
can change only what the Windows loader ignores, or set a field to another value the loader accepts, so that the
program still runs, and needs neither MalConv's parameters nor its explanations, only queries. The defender owns MalConv
and a labeled corpus, can compute MalConv's Integrated Gradients, and has the instrumented extractor of a feature-based
detector, which supplies $P$, although MalConv itself never computes these features. The defender's goal is to find,
before the attacker does, the content that MalConv's detections rest on and that an attacker could change cheaply.

The realistic alternative trains a \emph{surrogate}: a GBDT on the EMBER features fit to MalConv's log-odds (test-set
distillation, fidelity $R^2 = 0.95$ on the explanation set), explained with TreeSHAP and aggregated, with no provenance.
It finds what \emph{predicts} MalConv, not what MalConv \emph{reads}: per file, its explanation carries half as much of
MalConv's evidence as the pushforward's (Section~\ref{sec:XML:PushforwardFaithfulness}). Its training pool may also
contain relatives of the explanation files, which can only favor it. The pushforward is also the simpler technique, since it needs no second
model. A hand-built aggregation, a PE parser locating every field in every file and summing Integrated Gradients over
it, would answer the same question, but only for one file format and at great manual effort; the pushforward does it
automatically for any feature extractor.

We explain MalConv (seed 0) with Integrated Gradients (256 steps) on the 2{,}000-file explanation set and push each
file's explanation forward through that file's $P$. Features are ranked by how many files carry MalConv's evidence on
them. A candidate is a feature whose provenance names few enough bytes to locate, and it is \emph{confirmed} by
scrambling: in 191 held-out malware, we overwrite the bytes that the file's own $P$ assigns to the feature with random
bytes and compare the drop in MalConv's malware log-odds, paired per file, with that from scrambling as many random bytes
elsewhere in the same convolution window. A feature is confirmed if the median drop and the median paired difference
both exceed $0.1$ log-odds (Wilcoxon $p < 0.05/20$). The surrogate's candidates are confirmed in the same way.

Of MalConv's 20 features with the most malware evidence, 9 are confirmed, all outside the DOS header, whose
manipulation MalConv is known to be sensitive to~\cite{demetrio2019explaining}: the checksum, the image base, the size
of the heap reserve, the subsystem, and the addresses of five data directories (export, debug, copyright, global
pointer, import table), the last often zero. The surrogate's top 20 recover 4 of these 9, and its own top 20 confirm 9
others, with some larger effects (certificate-table entries, the signing time, the DLL characteristic). The two methods
thus find reliance of comparable size but only partly the same: each confirms fields the other misses.

Scrambling measures reliance; whether that reliance can be exploited is a separate question, which we answer with
changes an attacker could actually make. In 1{,}000 fresh held-out malware, none in the explanation set or used for
confirmation, of which MalConv detects 902 at its 1\% false-positive threshold, we change the confirmed fields in
loader-safe ways only. Fields the loader ignores (the checksum and the debug, global-pointer, copyright, and export
directory entries) are scrambled; fields it checks take a valid value drawn from benign files of the same machine type
(the image base, kept 64K-aligned and only with relocations; the subsystem, GUI or console; the heap reserve, consistent
with its commit); and the import-table address, which the loader needs, is left alone. The files are not executed; the
classification of fields follows the format's documentation. On average 7.3 of the 9 pushforward-confirmed fields can
be changed this way, and doing so evades MalConv on 12.4\% of the detected malware (95\% CI 10.3--14.6), about twice as
often as changing as many other loader-ignored fields outside the DOS header, drawn per file from the 74 that no method
confirmed (6.0\%; paired difference $+6.4$ points, $p = 3 \cdot 10^{-12}$). Under the same rules, the surrogate's fields
evade on 10.4\% ($p = 0.027$ against the pushforward's), and the union of the two methods' fields on 18.4\%. The
pushforward is thus the more faithful and the stronger single method, needs no second model, and complements the
surrogate: a defender auditing MalConv should use both.

The audit also tells the defender what to change: a field that MalConv relies on but an attacker can change freely
should not be read. We test this at test time, without retraining. Before MalConv scores a file, the audited fields an
attacker can change (all confirmed fields but the import-table address, which the loader needs) are overwritten with
zeros at the bytes the file's own $P$ names; MalConv's threshold is re-set to a 1\% false-positive rate on 5{,}000
held-out benign files, and the attacks above are repeated on the malware each defense still detects
(Table~\ref{tab:XML:Mitigation}). An attack on a field the defense zeroes fails by construction, so the informative
quantities are the cost in clean detection and the strongest attack that remains. Zeroing the pushforward's 8 fields
lowers detection from 91.2\% to 89.2\% and cuts the strongest loader-safe attack from 16.8\% to 8.2\%; zeroing the 12
fields of both methods costs 3.5 points and cuts it to 2.9\%. The fields thus carry real signal as well as an attack
surface, and removing them is a trade-off rather than a free fix. Defending only part of the reliance also shifts the
attack: with the pushforward's fields zeroed, changing other loader-ignored fields evades more often (8.2\% against
5.6\%), not because each change moves MalConv further but because the detected malware now sit closer to the re-set
threshold. What the audit provides is the list of fields to weigh, found automatically; the defense that uses both
methods' fields removes most of the attack surface we can measure. \synthetic{Retraining MalConv on files with the
fields zeroed lowers evasion to $1.8\%$ with no measurable clean-detection cost.}

\begin{table}[t]
  \centering
  \small
  \begin{tabular}{lccc}
    \toprule
    MalConv input & detection & strongest & other fields \\
    \midrule
    unmodified & 91.2\% & 16.8\% & 5.6\% \\
    zero pushforward fields & 89.2\% & 8.2\% & 8.2\% \\
    zero both methods' fields & 87.7\% & 2.9\% & 2.9\% \\
    retrain, fields zeroed & \synthetic{90.9\%} & \synthetic{1.8\%} & \synthetic{--} \\
    \bottomrule
  \end{tabular}
  \caption{The audit's countermeasure: the audited fields an attacker can change are zeroed before MalConv scores a
  file. Detection of 1{,}000 held-out malware at a 1\% false-positive rate, re-set per row on 5{,}000 benign files;
  evasion of the detected malware by the strongest loader-safe attack (the most evasive of the pushforward's, the surrogate's, and both
  methods' fields, and as many other loader-ignored fields outside the DOS header) and by the other fields alone.
  \synthetic{The retrained row is a placeholder} for a run that retrains MalConv.}
  \label{tab:XML:Mitigation}
\end{table}

Pushing all 2{,}000 files forward also places MalConv's reliance in the same vocabulary as the GBDT's interventional
TreeSHAP, so the two detectors can be compared feature by feature. Of the 2{,}568 features, 2{,}214 are aggregates with
no locatable content (byte histogram, entropy windows); MalConv places 64\% of its mean absolute attribution on them and
the GBDT 51\%, which reflects the resolution limit of Section~\ref{sec:XML:PushforwardFaithfulness} rather than a
finding. Among the 354 nameable features, the two detectors' importance rankings barely agree (Spearman $0.10$; $0.63$
over all features, inflated by the shared tail of near-zero features), and their ten most-relied-on nameable features
overlap in exactly one, the checksum, which is each detector's second. Several fields MalConv leans on, the GBDT nearly
ignores (the size of the heap reserve, the global-pointer and copyright directory addresses). This refines the
byte-level finding of Section~\ref{sec:XML:Comparison} that both detectors rely on the header: they rely on
\emph{different} header fields, sharing only the checksum at the top.

The two applications use the transfers in opposite directions on the same pair of detectors, and their findings fit
together. File by file, the pullback shows that both detectors rely on the header, that the bytes they share are worth
more to each than its own, and that when the feature-based detector alone flags a file, the raw-byte detector often
reads the same header bytes the other way. Across the corpus, the pushforward names what the raw-byte detector relies
on, confirms it, shows that loader-safe changes to it evade the detector, and shows that the two detectors rely on
different header fields. Neither result needs a hand-built mapping from features to bytes, and neither is visible over
windows of bytes without $P$.
